% !TeX root = ../main.tex

\chapter{实验结果与分析}
\label{ch:experiment}

\section{实验环境与设置}

本文实验围绕 LEVIR-Ship 数据集展开，采用统一的数据划分和评测口径，对 DRENet、FCOS 与 YOLO26 三类代表模型进行对比分析。核心评价指标包括 AP50、AP50-95、Precision、Recall 与 F1；为辅助分析工程部署能力，还统计了 FPS、参数量与 FLOPs 等效率指标。实验过程中，训练配置、权重文件、运行日志和结果文档均进行了统一留痕，以保证论文表格和系统演示能够回溯到具体实验产物。

需要特别说明的是，DRENet 与 YOLO26 采用统一的 512 输入主线，而 FCOS 当前可用的正式结果来自 MMDetection 默认训练口径，其输入策略为 1333\(\times\)800。基于这一客观差异，本文对三模型的使用口径作出如下区分：DRENet 与 YOLO26 作为论文主结论模型进行重点比较；FCOS 作为 anchor-free 参考基线参与对比，用于补充不同检测范式的讨论，而不用于强化“同输入尺度严格公平”的结论。

从工程角度看，本文不仅保留了训练与测试结果，也完成了系统侧的真实推理联调。这意味着本章中的定量分析并非孤立存在，而是能够与第五章中的系统实现相互印证。

\section{三模型主对比结果}

目前已完成并可追溯的三模型核心结果如表~\ref{tab:ch4_baselines} 所示。

\begin{table}[htbp]
  \centering
  \caption{三模型在 LEVIR-Ship 数据集上的主对比结果}
  \label{tab:ch4_baselines}
  \begin{tabular}{lcccccc}
    \toprule
    模型 & AP50 & AP50-95 & Precision & Recall & F1 & 说明 \\
    \midrule
    DRENet & 0.7949 & 0.2919 & 0.4927 & 0.8511 & 0.6241 & 论文主结论模型 \\
    FCOS & 0.7700 & 0.2850 & -- & 0.4050$^{*}$ & -- & anchor-free 参考基线 \\
    YOLO26 & 0.7950 & 0.3170 & 0.8430 & 0.7220 & 0.7778 & 论文主结论模型 \\
    \bottomrule
  \end{tabular}
  \\
  \parbox{0.92\textwidth}{\footnotesize 注：FCOS 当前可直接回填的 0.4050 来自正式 run 中的 \texttt{AR@100} 指标，本文仅用其补充说明检出能力与参考表现，不与 DRENet 和 YOLO26 的 Recall 做严格等值比较。}
\end{table}

表中可以看出，YOLO26 与 DRENet 在 AP50 上表现接近，分别为 0.7950 和 0.7949，均明显优于 FCOS 的 0.7700。若仅从主指标 AP50 观察，DRENet 与 YOLO26 的整体检测能力相当。但从其他指标进一步分析，可以发现二者的侧重点并不相同。

首先，从 AP50-95 来看，YOLO26 达到 0.3170，高于 DRENet 的 0.2919 和 FCOS 的 0.2850，说明 YOLO26 在更严格的 IoU 阈值下具有更好的定位质量。其次，从 Precision 与 Recall 的平衡关系看，DRENet 的 Recall 达到 0.8511，显著高于 YOLO26 的 0.7220，说明其更倾向于检出尽可能多的目标；而 YOLO26 的 Precision 达到 0.8430，明显高于 DRENet 的 0.4927，说明其误检控制能力更强。因此，二者虽然在 AP50 上接近，但在“召回优先”与“精度优先”的取舍上具有明显差异。

FCOS 的 AP50 与 AP50-95 均低于 DRENet 和 YOLO26，这说明在当前实验设置下，anchor-free 通用基线在遥感微小船舶场景中的整体表现尚不占优。但 FCOS 仍然具有参考意义：它为本文提供了与专项方法、one-stage 方法不同的检测范式参照，也为后续分析“是否必须采用专项设计”提供了比较基础。

\section{效率与部署属性分析}

除精度指标外，工程部署场景还需要关注模型在速度与复杂度方面的表现。当前可用效率结果如表~\ref{tab:ch4_efficiency} 所示。

\begin{table}[htbp]
  \centering
  \caption{三模型效率与复杂度对比}
  \label{tab:ch4_efficiency}
  \begin{tabular}{lccccl}
    \toprule
    模型 & FPS & Params(M) & FLOPs(G) & 输入尺寸 & 说明 \\
    \midrule
    DRENet & 121.8993 & 4.7882 & 4.2057 & 512 & 同口径效率测试 \\
    FCOS & 45.9559 & 32.1664 & 123.0 & 1333\(\times\)800 & 输入策略不同，仅作参考 \\
    YOLO26 & 69.0625 & 2.5042 & 3.6939 & 512 & 同口径效率测试 \\
    \bottomrule
  \end{tabular}
\end{table}

结合表~\ref{tab:ch4_efficiency} 可以看到，DRENet 在当前记录中取得了最高 FPS，说明其在本文所采用的环境和测量方式下具有较好的运行速度；YOLO26 的参数量和 FLOPs 更低，体现出较强的轻量化优势；FCOS 由于输入策略和实现框架差异，其效率数据更多反映“当前历史正式 run 的参考成本”，不宜与前两者进行严格同口径比较。

从部署角度看，这些结果进一步支撑了前文的判断：如果更关注漏检控制，DRENet 是值得优先考虑的候选；如果更关注轻量部署、整体精度和系统接入便利性，YOLO26 更具优势。FCOS 则可以作为补充模型保留在系统与论文中，用于完善范式比较，但不宜作为系统默认主模型。

\section{模型逐项分析}

\subsection{DRENet 结果分析}

DRENet 是面向遥感微小船舶任务的专项方法，其最突出特点体现在较高的 Recall 指标上。当前实验中，DRENet 的 Recall 达到 0.8511，明显高于 YOLO26，这说明 DRENet 在复杂背景与弱纹理目标场景下更倾向于给出候选框，从而减少漏检。对于海上监管、目标排查等更重视“尽可能发现目标”的任务，这种特性具有实际价值。

但与此同时，DRENet 的 Precision 为 0.4927，低于 YOLO26。这表明其在提升召回的同时也引入了较多误检，特别是在近岸背景、港口设施和纹理复杂区域中更容易受到干扰。因此，DRENet 更适合作为强调召回优先的主方法，若要进一步用于实际部署，还需要结合阈值调优、后处理或多模型融合来抑制误检。

\subsection{FCOS 结果分析}

FCOS 作为本文选取的 anchor-free 代表模型，其 AP50 为 0.7700，整体结果弱于 DRENet 与 YOLO26。从最终表现看，通用 anchor-free 路线在当前中分辨率微小船舶场景中的优势并不明显，至少在本文现有实验条件下尚未体现出超越专项方法或成熟 one-stage 模型的能力。

不过，FCOS 的意义并不在于追求最佳结果，而在于提供一个结构清晰、范式明确的参考基线。通过引入 FCOS，本文能够更完整地回答“专项设计是否必要”“通用 anchor-free 方法在该任务中能达到怎样的水平”等问题。结合训练记录看，FCOS 也为后续更规范的框架复现、输入策略分析和稳定性讨论保留了空间。

\subsection{YOLO26 结果分析}

YOLO26 在当前实验中取得了最优的 AP50-95、Precision 和 F1，同时 AP50 与 DRENet 基本持平，说明其在检测精度、定位质量和误检控制之间表现出较好的综合平衡。尤其是 Precision 达到 0.8430，说明在当前阈值设置下，YOLO26 对复杂背景下的伪目标抑制能力较强。

从工程角度看，YOLO26 还具有训练链路成熟、部署接口清晰和系统接入成本较低等优势。因此，若从“本科毕设需要稳定可演示系统”的角度考虑，YOLO26 是很有竞争力的默认模型候选。与 DRENet 结合时，YOLO26 可以在整体精度和误检控制方面提供互补。

\section{定性样例与误差分析}

定量结果能够反映整体性能，但难以完整解释模型在复杂场景中的具体行为。因此，本文进一步结合系统导出的可视化结果，对成功检测、漏检与误检样例进行分析。当前样例已按“模型 $\rightarrow$ 场景类型”归档，能够直接服务论文插图与答辩展示。

\subsection{成功检测样例}

在成功检测样例中，DRENet 往往能够在海面纹理较弱、目标尺寸较小的情况下给出有效响应，说明其对弱目标具有一定增强作用。YOLO26 在多目标和常规场景中的整体输出较为整洁，预测框与标注框匹配较好。FCOS 在部分样例中也能够给出合理定位，但整体稳定性略弱于前两者。

\begin{figure}[htbp]
  \centering
  \fbox{\parbox{0.88\textwidth}{
  图占位：成功检测样例图组。建议按 DRENet、YOLO26、FCOS 顺序排布，每组展示同类成功样例各 1 张，并在图注中说明图像来源于 qualitative 目录。
  }}
  \caption{三模型成功检测样例对比}
  \label{fig:ch4_success_cases}
\end{figure}

\subsection{漏检样例}

在漏检样例中，YOLO26 更容易在局部弱目标、对比度较低或被背景纹理淹没的场景下丢失部分目标；DRENet 虽然整体召回较高，但在极度稀疏或局部低对比区域仍可能出现覆盖不完整的情况。FCOS 的漏检现象相对更明显，这与其当前主结果低于 DRENet 和 YOLO26 的定量结论保持一致。

\begin{figure}[htbp]
  \centering
  \fbox{\parbox{0.88\textwidth}{
  图占位：漏检样例图组。建议包含 DRENet miss、YOLO26 miss 和 FCOS miss 三组对照，并在图注中说明标签图与预测图的对应关系。
  }}
  \caption{三模型漏检样例对比}
  \label{fig:ch4_miss_cases}
\end{figure}

\subsection{误检样例}

在误检样例中，DRENet 更容易将近岸纹理、港口设施或海面复杂波纹误识别为船舶目标，这与其高 Recall、低 Precision 的定量特征一致。YOLO26 在误检抑制方面整体更稳，但在纹理复杂区域仍可能出现额外框或重复框。FCOS 也存在一定误检现象，其 global-best 与 stable 两套样例可分别对应论文口径和系统口径展示。

\begin{figure}[htbp]
  \centering
  \fbox{\parbox{0.88\textwidth}{
  图占位：误检样例图组。建议展示 DRENet false positive、YOLO26 false positive 和 FCOS false positive，对应 qualitative 目录中的现有样例。
  }}
  \caption{三模型误检样例对比}
  \label{fig:ch4_false_positive_cases}
\end{figure}

综合定性样例可以看出，DRENet 与 YOLO26 的差异并不只是体现在数值上，更体现在错误类型分布上。DRENet 更强调“先检出来”，因此在复杂背景下更容易增加误检；YOLO26 更强调“框更稳、输出更干净”，但在弱目标场景下可能牺牲部分召回。FCOS 则为这一比较提供了 anchor-free 路线的补充参考。

\section{系统口径与论文口径说明}

为了保证论文论证与系统实现之间的一致性，本文对模型使用口径进行了专门区分。在论文主结论中，重点使用 DRENet 与 YOLO26 的正式结果，并将 FCOS 保留为参考基线；在系统实现中，除了单模型推理外，还提供融合模式和默认权重配置，用于增强系统演示时的结果完整性。对于 FCOS，还进一步区分了 \textit{global-best} 与 \textit{stable-best} 两类权重：前者更适合论文中展示其最佳离线结果，后者更适合作为系统联调中的稳定候选。

这种区分并不意味着论文和系统相互割裂，而是强调二者服务于不同目标：论文更关注实验结论的解释力和口径一致性，系统更关注任务执行的稳定性与展示效果。明确这一点，有助于在答辩中更清楚地解释模型选择与结果展示策略。

\section{阶段性不足与后续补充方向}

尽管本文已经形成较完整的主对比结果和定性样例，但实验部分仍存在若干阶段性不足。首先，当前消融实验尚未形成完整离线表格，阈值与分辨率的进一步对照仍需要补充。其次，FCOS 的效率结果由于输入策略不同，只能作为参考值，而不能与 DRENet 和 YOLO26 做完全同口径比较。再次，系统中的融合模式已经能够稳定运行，但目前尚未形成统一离线 AP 指标表，因此本文暂不把融合结果作为主结论展开论证，而是主要将其用于系统展示和结果互补分析。

对于本科毕业设计而言，以上不足并不会破坏整篇论文的完整性，因为主结论已经由 DRENet 与 YOLO26 的正式结果支撑，FCOS 的角色也已明确界定。后续若有进一步时间和算力条件，可围绕阈值消融、融合离线评估和更严格统一口径效率测试继续完善。

\section{本章小结}

本章在统一评测协议下，对 DRENet、FCOS 和 YOLO26 三种模型进行了定量与定性分析。结果表明，DRENet 在召回方面表现突出，更适合强调漏检控制的场景；YOLO26 在 AP50、AP50-95、Precision 和 F1 等指标上表现均衡，更适合作为轻量部署与系统默认模型候选；FCOS 则提供了 anchor-free 路线的参考基线。结合效率分析、误差样例和系统口径说明，可以认为本文已经形成较完整的实验论证框架，为下一章的系统设计与实现提供了直接依据。

\section{消融实验}

为进一步增强主结论的可解释性，本文补充两组离线消融分析。第一组为三模型统一的置信度阈值消融，即在固定 IoU 阈值下调节 conf\_threshold，比较 DRENet、FCOS 和 YOLO26 在 Precision 与 Recall 平衡上的变化趋势；第二组为输入尺寸敏感性分析，即在不重新训练模型的前提下，仅改变推理阶段输入尺寸，观察模型对小目标尺度放大的响应差异。需要说明的是，输入尺寸组属于推理阶段敏感性验证，不等同于重新训练后的完整尺寸消融，其主要作用是解释不同模型在现有正式权重基础上对输入尺度扰动的行为差异。

在具体实施上，三模型统一阈值消融能够直接形成横向可比结果；而输入尺寸敏感性分析则需要遵守各模型实现边界。YOLO26 与 FCOS 支持推理阶段输入尺寸覆写，因此可以开展 512、640、800 等多点对比；DRENet 当前本地插件在非 512 输入下会触发结构 shape mismatch，因此不纳入统一尺寸主表，而是将该现象作为实现限制单独记录。对于 FCOS，还需额外说明其正式参考基线来自历史默认训练口径，而非本文后续尝试过的统一 512 口径，因此尺寸敏感性结果只用于行为分析，不反向改写主对比结论。

从分析框架看，若某一模型在更高推理尺寸下 Recall 提升更明显，可说明该模型对微小目标可见性的依赖更强；若某一模型在 conf 提高后 Precision 上升而 Recall 快速下降，则说明其检测输出更依赖后处理阈值进行筛选。通过阈值扰动与尺寸扰动两条分析线，本文不仅关注最终指标高低，也进一步讨论不同检测范式在参数变化下的行为模式差异，从而为模型选择与系统默认参数设定提供依据。
